\section{Conclusion and Future Work}\label{sec:concl}

This project asked whether a physics-informed, mask-aware recurrent network predicts hourly HVAC energy more accurately than conventional models, including the published baseline, when a building's sensors are incomplete. On LBNL Building 59, with the baseline's own target, split and metrics, the PI-GRU beats the published BNN of \citet{mahajan2024bnn} on every metric (RMSE 9.01 against 9.65\,kWh, MAE 5.70 against 7.06\,kWh, MAPE 0.22 against 0.35) and stays 20\% better than the re-implemented baseline under 40\% contiguous sensor outages. XGBoost on the same mask-aware features is more accurate still in RMSE (7.97\,kWh).

The objectives were met. A reproducible benchmark with identical seeded masks was built (O1); the baseline was re-implemented and compared with five conventional models and two ablations (O2); accuracy, robustness, physical consistency and cost were evaluated with five seeds, bootstrap intervals and Holm-adjusted tests (O3); and the baseline was beaten on its own protocol (O4). The hypotheses were answered honestly: the PI-GRU did not beat the best conventional model (H1 rejected), the physics term did not improve accuracy over its matched twin (H2 not supported), and it lowered the heat-balance residual on complete data without a meaningful accuracy cost (H3 partly supported). Random sensor loss proved almost harmless and contiguous outages harmful, and mask awareness, not physics, gave the largest robustness gain.

\textbf{Future work.}
\begin{itemize}[nosep]
\item Replicate on buildings where the metered energy covers the whole plant and where indoor temperatures vary more, so that the RC parameters are identifiable.
\item Use a second-order RC model with zone temperatures, or a hard-constrained architecture such as a physically consistent network \citep{dinatale2022pcnn}, where zone-level data allow it.
\item Add realistic sensor faults (drift, stuck values, bias) to the outage generator.
\item Combine the strong tabular model with the recurrent model, for example by stacking, and add calibrated uncertainty under sensor loss, which was the focus of the baseline.
\item Evaluate transfer across buildings from the same data programme.
\end{itemize}
